
\documentclass{article}
\usepackage{colortbl}
\usepackage{makecell}
\usepackage{multirow}
\usepackage{supertabular}

\begin{document}

\newcounter{utterance}

\centering \large Interaction Transcript for game `cladder', experiment `full\_v1.5\_default', episode 1796 with refineonly{-}llama{-}3.1{-}8b{-}instruct.
\vspace{24pt}

{ \footnotesize  \setcounter{utterance}{1}
\setlength{\tabcolsep}{0pt}
\begin{supertabular}{c@{$\;$}|p{.15\linewidth}@{}p{.15\linewidth}p{.15\linewidth}p{.15\linewidth}p{.15\linewidth}p{.15\linewidth}}
    \# & \multicolumn{2}{c}{Player} && \multicolumn{2}{c}{Game Master} \\
    \hline

    \theutterance \stepcounter{utterance}  
    & & & \multicolumn{4}{p{0.6\linewidth}}{
        \cellcolor[rgb]{0.9,0.9,0.9}{
            \makecell[{{p{\linewidth}}}]{
                \texttt{\tiny{[P1$\langle$GM]}}
                \texttt{You are an expert in causal inference. The following question is not a typical commonsense query, but rather a meticulously designed question created by a professor specializing in causal inference, intended to assess the students' mastery of the course content.} \\
\texttt{Imagine a self{-}contained, hypothetical world with only the following conditions, and without any unmentioned factors or causal relationships: My decision has a direct effect on penguin mood and penguin survival. Penguin mood has a direct effect on penguin survival. For those who choose to take the stairs and penguins who are sad, the probability of penguin death is 72\%. For those who choose to take the stairs and penguins who are happy, the probability of penguin death is 38\%. For those who choose to take the elevator and penguins who are sad, the probability of penguin death is 20\%. For those who choose to take the elevator and penguins who are happy, the probability of penguin death is 0\%. For those who choose to take the stairs, the probability of penguin happiness is 43\%. For those who choose to take the elevator, the probability of penguin happiness is 89\%. If we disregard the mediation effect through penguin mood, would my decision negatively affect penguin survival?} \\
\texttt{Start your answer with 'Yes' or 'No', followed by additional reasoning or evidence to support your explanation.} \\
            }
        }
    }
     \\ \\

    \theutterance \stepcounter{utterance}  
    & & \multicolumn{4}{p{0.6\linewidth}}{
        \cellcolor[rgb]{0.95,0.95,0.95}{
            \makecell[{{p{\linewidth}}}]{
                \texttt{\tiny{[P1$\leftrightarrow$P]}}
                \texttt{{[}DRAFT{]}} \\
\texttt{No.} \\
\\ 
\texttt{To assess the effect of my decision on penguin survival, we can use the potential outcomes framework. Let's define two potential outcomes: penguin survival when taking the stairs (S) and penguin survival when taking the elevator (E). We can then calculate the average treatment effect (ATE) as the difference between these two potential outcomes.} \\
\\ 
\texttt{First, we need to calculate the probability of penguin survival for each scenario:} \\
\\ 
\texttt{{-} For those who choose to take the stairs and penguins who are sad, the probability of penguin death is 72\%, so the probability of penguin survival is 28\%.} \\
\texttt{{-} For those who choose to take the stairs and penguins who are happy, the probability of penguin death is 38\%, so the probability of penguin survival is 62\%.} \\
\texttt{{-} For those who choose to take the elevator and penguins who are sad, the probability of penguin death is 20\%, so the probability of penguin survival is 80\%.} \\
\texttt{{-} For those who choose to take the elevator and penguins who are happy, the probability of penguin death is 0\%, so the probability of penguin survival is 100\%.} \\
\\ 
\texttt{We can then calculate the probability of penguin happiness for each scenario:} \\
\\ 
\texttt{{-} For those who choose to take the stairs, the probability of penguin happiness is 43\%.} \\
\texttt{{-} For those who choose to take the elevator, the probability of penguin happiness is 89\%.} \\
\\ 
\texttt{Using the law of total probability, we can calculate the overall probability of penguin survival for each decision:} \\
\\ 
\texttt{{-} For those who choose to take the stairs, the overall probability of penguin survival is:} \\
\texttt{  (0.43 * 0.62) + (0.57 * 0.28) = 0.373 + 0.158 = 0.531} \\
\texttt{{-} For those who choose to take the elevator, the overall probability of penguin survival is:} \\
\texttt{  (0.43 * 1) + (0.57 * 0.8) = 0.43 + 0.456 = 0.886} \\
\\ 
\texttt{The average treatment effect (ATE) is the difference between these two potential outcomes:} \\
\\ 
\texttt{ATE = 0.886 {-} 0.531 = 0.355} \\
\\ 
\texttt{Since the ATE is positive, my decision to take the elevator would lead to a higher probability of penguin survival compared to taking the stairs. Therefore, disregarding the mediation effect through penguin mood, my decision would not negatively affect penguin survival.} \\
            }
        }
    }
    & \\ \\

    \theutterance \stepcounter{utterance}  
    & \multicolumn{4}{p{0.6\linewidth}}{
        \cellcolor[rgb]{0.561,0.737,0.561}{
            \makecell[{{p{\linewidth}}}]{
                \texttt{\tiny{[P1$\rangle$GM]}}
                \texttt{No.} \\
\\ 
\texttt{The average treatment effect (ATE) is the difference between the overall probability of penguin survival when taking the stairs and when taking the elevator. } \\
\\ 
\texttt{ATE = 0.886 {-} 0.531 = 0.355} \\
\\ 
\texttt{Since the ATE is positive, my decision to take the elevator would lead to a higher probability of penguin survival compared to taking the stairs. Therefore, disregarding the mediation effect through penguin mood, my decision would not negatively affect penguin survival.} \\
            }
        }
    }
    & & \\ \\

    \theutterance \stepcounter{utterance}  
    & & & \multicolumn{2}{p{0.3\linewidth}}{
        \cellcolor[rgb]{0.9,0.9,0.9}{
            \makecell[{{p{\linewidth}}}]{
                \texttt{\tiny{[GM$|$GM]}}
                \texttt{No} \\
            }
        }
    }
    & & \\ \\

    \theutterance \stepcounter{utterance}  
    & & & \multicolumn{2}{p{0.3\linewidth}}{
        \cellcolor[rgb]{0.9,0.9,0.9}{
            \makecell[{{p{\linewidth}}}]{
                \texttt{\tiny{[GM$|$GM]}}
                \texttt{yes} \\
            }
        }
    }
    & & \\ \\

    \theutterance \stepcounter{utterance}  
    & & & \multicolumn{2}{p{0.3\linewidth}}{
        \cellcolor[rgb]{0.9,0.9,0.9}{
            \makecell[{{p{\linewidth}}}]{
                \texttt{\tiny{[GM$|$GM]}}
                \texttt{game\_result = LOSE} \\
            }
        }
    }
    & & \\ \\

\end{supertabular}
}

\end{document}
